\selectlanguage{russian}
\section{Обзор и постановка задачи}
\label{sec:Background} \index{Background}

\subsection{Обозначения}

Основные обозначения приведены в Таблице~\ref{tab:notation}. На общем уровне
$p_\theta$ обозначает условное распределение модели, а $q_\phi$ --- вариационное
распределение. Их конкретные параметризации вводятся в Разделе~\ref{sec:Method}.

\begin{table}[H]
\centering
\caption{Обозначения, используемые в работе.}
\label{tab:notation}
\begin{tabular}{ll}
\toprule
Символ & Значение \\
\midrule
$p^*$ & истинное распределение данных \\
$p_\theta(y\mid x)$ & условное распределение модели с параметрами $\theta$ \\
$q_\phi$ & вариационное распределение с параметрами $\phi$ \\
$p_Z$ & априорное распределение скрытой переменной \\
$N$ & число независимых групп \\
$M$ & число точек в группе \\
$x^n_m,\,y^n_m$ & $m$-е значения $x$ и $y$ в группе $n$ \\
$x^n,\,y^n$ & векторы $(x^n_1,\dots,x^n_M)$ и $(y^n_1,\dots,y^n_M)$ \\
$\sigma^n\in S_M$ & скрытая внутригрупповая перестановка группы $n$ \\
$S_M$ & симметрическая группа перестановок множества $\{1,\dots,M\}$ \\
$z$ & мягкая матрица назначений из $[0,1]^{M\times M}$ \\
$\mathcal{B}_M$ & многогранник дважды стохастических матриц \\
$\mathbb{S}^1=\mathbb{R}/\mathbb{Z}$ & одномерный тор \\
$\tilde{\mathcal{N}}$ & обёрнутое нормальное распределение на торе \\
$\mathrm{KL}(\cdot\,\|\,\cdot)$ & дивергенция Кульбака--Лейблера \\
$K$ & число компонент смеси в конкретной параметризации \\
$\tau$ & температура Sinkhorn-преобразования \\
$\beta$ & вес KL-регуляризатора \\
$\epsilon$ & стандартное отклонение порождающего шума \\
\bottomrule
\end{tabular}
\end{table}

\subsection{Формальная постановка задачи}

Даны $N$ независимых групп. Для группы $n\in\{1,\dots,N\}$ наблюдаются два набора по $M$
элементов,
\begin{equation}
x^n=(x^n_1,\dots,x^n_M), \qquad
y^n=(y^n_1,\dots,y^n_M).
\end{equation}
До потери соответствия элементы были получены как парные наблюдения из неизвестного
совместного распределения. Затем значения $y$ были переупорядочены скрытой перестановкой
$\sigma^n\in S_M$. Следовательно, $x^n_m$ образует истинную пару с
$y^n_{\sigma^n(m)}$, но значение $\sigma^n$ при обучении неизвестно.

Требуется оценить условное распределение $p^*(y\mid x)$ с помощью параметрической модели
$p_\theta(y\mid x)$. При условной независимости пар внутри группы и скрытой перестановке
групповое правдоподобие имеет вид
\begin{align}
p_\theta(y^n\mid x^n)
&=\sum_{\sigma\in S_M}p(\sigma)
  p_\theta(y^n\mid x^n,\sigma) \nonumber\\
&=\sum_{\sigma\in S_M}p(\sigma)
  \prod_{m=1}^{M}p_\theta\bigl(y^n_{\sigma(m)}\mid x^n_m\bigr).
\label{eq:global-marginal}
\end{align}
При отсутствии априорной информации о перестановках используется равномерное распределение
\begin{equation}
p(\sigma)=\frac{1}{M!}, \qquad \sigma\in S_M.
\end{equation}
Целью обучения служит максимизация суммы логарифмов группового правдоподобия по параметрам
$\theta$. Точное вычисление~\eqref{eq:global-marginal} требует учёта $M!$ перестановок и
становится трудновычислимым при росте $M$.

\subsection{Аппроксимация Байера и др.}

Постановка восстановления распределения по группам непарных данных предложена Байером и
др.~\cite{beier2024transfer}. Авторы начинают с точного правдоподобия, соответствующего
усреднению произведения попарных плотностей по всем перестановкам, как
в~\eqref{eq:global-marginal}. Получаемая оценка максимального правдоподобия имеет
факториальную вычислительную сложность по размеру группы.

Для получения вычислимого функционала Байер и др. заменяют зависимые координаты случайно
переставленной группы их независимыми маргинальными распределениями. В обозначениях
настоящей работы приближённое отрицательное логарифмическое правдоподобие записывается как
\begin{equation}
\widetilde{\mathcal{J}}_N(\theta)
=-\frac{1}{NM}\sum_{n=1}^{N}\sum_{j=1}^{M}
\log\left[
\frac{1}{M}\sum_{k=1}^{M}p_\theta\bigl(y^n_j\mid x^n_k\bigr)
\right]
\;\longrightarrow\;\min_\theta.
\label{eq:beier-approximation}
\end{equation}
Вместо суммы по $M!$ перестановкам здесь для каждой группы вычисляются $M^2$ попарных
плотностей. Цена аппроксимации состоит в отказе от совместного распределения назначений:
выражение~\eqref{eq:beier-approximation} не задаёт единую биекцию между элементами группы.

В работе Байера и др. пространство допустимых плотностей строится с помощью энтропийных
транспортных ядер. Параметры плотности оцениваются минимизацией приближённого отрицательного
логарифмического правдоподобия; после дискретизации для этого используется модифицированный
EMML-алгоритм с ограничениями на переходные вероятности. Таким образом, транспортное ядро
задаёт семейство моделей, а обучение остаётся оцениванием методом максимального
правдоподобия.

\subsection{Жёсткое назначение в нейросетевых моделях}

Deep Perm-Set Net~\cite{rezatofighi2018perm} и DETR~\cite{carion2020detr} используют
Hungarian matching для выбора дискретного соответствия между предсказанным и целевым
наборами. После выбора соответствия параметры модели обновляются по функции потерь для
сопоставленных элементов; дифференцирование через сам алгоритм назначения этим методам не
требуется.

Такое жёсткое назначение нельзя непосредственно использовать как репараметризуемый
латентный слой вариационного автоэнкодера: Hungarian matching является дискретной
недифференцируемой операцией и не передаёт градиент от вариационной целевой функции к
параметрам распределения назначений. Непрерывная Sinkhorn-релаксация сохраняет мягкие веса
возможных назначений и допускает сквозное дифференцирование. Благодаря этому нейросетевую
версию модели можно обучать по мини-батчам стохастическим градиентным методом, а обученный
энкодер амортизирует вывод для новых групп.
